% Section 1. Placeholder: waits on the owner's PhD thesis (plan.md, section 1).
\section{Introduction}
\label{sec:intro}

\placeholder{Section~1 waits on the owner's PhD thesis (Linhares, \emph{Industrial
pattern sequencing problems: some complexity results and new local search
models}, 2001), which is not yet held. It is to record how the twelve problems of
Table~1 of \citet{LY2002} were collected, what the thesis proved and what it
only cited, and why the table appeared with no proofs and no graph
(\texttt{paper2/plan.md}, section~1). Kashiwabara and Fujisawa (1979), Table~1's
interval-thickness reference, is still missing.}

\paragraph{The pathwidth complex.} \citet{LY2002} collected, in their
Table~1, twelve problems from operations research, VLSI design and graph theory
that had been studied independently and that compute, on corresponding inputs,
``a function $f(\Pi)$ that is either equal to the number of open stacks or
closely related to it (plus or minus one)'' \citep[p.~1764]{LY2002}. The table
gave references, but no proofs and no graph. Table~\ref{tab:ly2002} reproduces
it as printed. We call this cluster \emph{the pathwidth complex}. Pathwidth is
both the name the literature uses most (Section~\ref{sec:names}) and the
quantity through which every relation in the table is proved
(Section~\ref{sec:equiv}). Of the twelve rows, eight are exactly pathwidth
plus a constant, two lie within a band of width one or two, and two are false
as Table~1 states them (Table~\ref{tab:master}).

\begin{table}[tbp]
\centering
\caption{Display of a family of equivalent problems studied independently on
the literature. Reproduced from \citet[Table~1, p.~1764]{LY2002}, with the
caption, row order and wording as printed; the references are the ones the
original cites, given here by author and year.}
\label{tab:ly2002}
\small
\begin{tabular}{@{}lll@{}}
\toprule
Problem & Discipline & References \\
\midrule
MOSP & Operations research & \citet{Yanasse1997b}; \citet{LY2002ref4} \\
Gate matrix layout & VLSI design & \citet{LY2002ref6}; \citet{LY2002ref8} \\
One-dimensional logic & VLSI design & \citet{LY2002ref7} \\
PLA folding & VLSI design & \citet{LY2002ref6} \\
Interval thickness & Graph theory & \citet{LY2002ref5} \\
Node search game & Graph theory & \citet{LY2002ref9} \\
Edge search game & Graph theory & \citet{LY2002ref10} \\
Narrowness & Graph theory & \citet{LY2002ref11} \\
Split bandwidth & Graph theory & \citet{LY2002ref12} \\
Graph path-width & Graph theory & \citet{LY2002ref13} \\
Edge separation & Graph theory & \citet{LY2002ref14} \\
Vertex separation & Graph theory & \citet{LY2002ref13} \\
\bottomrule
\end{tabular}
\end{table}

\paragraph{Contributions (draft list, for the introduction to argue).}
\begin{itemize}
\item A census of how much each of the twelve names is used, after removing the
  search hits that use the words in other senses; pathwidth has more relevant
  works than the other eleven names together, three times over
  (Section~\ref{sec:names}).
\item Each problem of Table~1 stated as an instance and a question, with the
  graph it lives on and its exact relation to pathwidth, every relation proved
  in Lean; two of Table~1's rows are false as stated (Section~\ref{sec:equiv}).
\item The exact open-stacks search of \citet{ChuStuckey2009} read as a
  vertex-separation search: two of its published dominance theorems are false,
  as are their restatements by \citet{Chu2011} and \citet{Fink2012}; the repair,
  its cost, and a Lean proof that the repaired search is sound; and a proof
  object for its refutations with an independent checker
  (Section~\ref{sec:search}).
\item A dataset of 17,714 isomorphism classes from twenty collections, 16,087 of
  them with a certified pathwidth and a witness layout, which through
  Section~\ref{sec:equiv} answers every problem in Table~1's exact core
  (Section~\ref{sec:dataset}).
\end{itemize}

\paragraph{Notation.} Graphs are finite and simple, $G=(V,E)$, $n=|V|$. A
\emph{layout} is a bijection $L:V\to\{1,\dots,n\}$. For a 0/1 matrix $M$ with
rows $R$ (customers, nets, piece types) and columns $C$ (patterns, gates,
products), the \emph{MOSP graph} $G_M$ has vertex set $R$, with $r\sim r'$ when
some column has a 1 in both rows; each column becomes a clique on its rows. For
an order $\pi$ of the columns, row $r$ is \emph{active} at position $j$ if it
has 1s in columns $x,y$ with $\pi(x)\le j\le\pi(y)$.
